% TOP_PROBLEM: {"id": 3271, "title": "Arc-disjoint directed cycles in regular directed graphs", "category_id": 3, "category": {"id": 3, "name": "graph_theory", "display_name": "Graph Theory", "description": "Problems involving graphs, networks, and their properties.", "slug": "graph-theory", "order_index": 3, "created_at": "2026-07-31T15:26:25.671Z"}, "status": "open", "problem_number": "OPG-49573", "set_id": 12}
% FIRST_POSTED: 2026-10-02
% ATTEMPT_STATUS: unresolved
% UnsolvedMath ID 3271; source catalogue code OPG-49573
% !TeX program = lualatex
% GPT-6 Astra Ultra research attempt; independently checked partial work, 2026-10-02.
\documentclass[11pt,a4paper]{article}
\usepackage[margin=25mm]{geometry}
\usepackage{fontspec}
\setmainfont{lmroman10-regular.otf}[BoldFont=lmroman10-bold.otf,
 ItalicFont=lmroman10-italic.otf,BoldItalicFont=lmroman10-bolditalic.otf]
\usepackage{amsmath,amssymb,mathtools}
\usepackage{unicode-math}
\setmathfont{latinmodern-math.otf}
\AtBeginDocument{\let\setminus\smallsetminus}
\usepackage{xurl,hyperref}
\hypersetup{colorlinks=true,urlcolor=blue}
\setcounter{secnumdepth}{0}
\setlength{\parindent}{0pt}
\setlength{\parskip}{5pt}
\setlength{\emergencystretch}{3em}
\begin{document}
\section*{Arc-disjoint directed cycles in regular directed graphs}\label{3271}
{\small UnsolvedMath ID 3271 · OPG-49573 · Graph Theory · OpenGarden}
\subsection{Definitions and mathematical statement}
Fix an integer $k\ge1$. A nonempty finite $k$-regular digraph $D=(V,A)$ here has
$d^+(v)=d^-(v)=k$ at every vertex. Assume it is loopless and has no
parallel arcs in the same direction. Opposite arcs are permitted and
are distinct members of $A$.

A directed cycle consists of distinct vertices joined cyclically by
arcs oriented in the direction of traversal. Its length is its number
of arcs; in particular an opposite pair of arcs is a cycle of length
two. A collection of cycles is arc-disjoint if no arc belongs to two
members. Vertices may be shared by many cycles.

The conjecture asks whether every such $D$ contains at least
\[
                             \binom{k+1}{2}=\frac{k(k+1)}2
\]
pairwise arc-disjoint directed cycles. It is enough to find exactly
that many, since extra cycles can be discarded. There is no requirement
that their union span all vertices or exhaust all arcs.

The coefficient is sharp for the complete symmetric digraph on $k+1$
vertices: it has exactly $k(k+1)$ arcs and every cycle uses at least two,
so an arc-disjoint collection contains at most $k(k+1)/2$ cycles. Taking
one digon for each unordered vertex pair attains that count. The
regularity convention is important: ``$k$-regular'' means both indegree
and outdegree equal $k$, not merely total degree $k$. The proposed lower
bound depends on this local degree, independently of the total order
of the digraph.
\subsection{Short English statement}
Must a digraph with indegree and outdegree k at every vertex contain at least k(k+1)/2 arc-disjoint directed cycles?
\subsection{Research attempt: Eulerian factorization and the cases $k=1,2$}
\paragraph{Scope and literature check.}
GPT-6 Astra Ultra, 2 October 2026. Regularity means both degrees equal
$k$, and opposite arcs form legitimate two-cycles. Alon,
McDiarmid and Molloy [S3] state the conjecture and prove substantially
stronger general lower bounds than the elementary linear estimates
below, including a bound of quadratic order. They also record the
conjecture for small degree parameters. The present proof is a
self-contained restricted verification, not a new best bound or a
claim that the cases established here were previously unknown.

\paragraph{1. Cycle decomposition of a balanced digraph.}
A finite digraph whose indegree equals its outdegree at every vertex
can have its entire arc set partitioned into directed simple cycles.
If an arc remains, start following arcs in its nonempty support.
A vertex reached by an incoming arc has a positive outgoing degree
as well, so the walk can continue until it repeats a vertex. Extract
a simple directed cycle and delete its arcs. Balance is preserved
because one entering and one leaving arc were removed at every
vertex of that cycle. Induction on the number of arcs completes the
decomposition. Isolated vertices cause no difficulty.

\paragraph{2. A spanning factorization gives a first lower bound.}
Make a bipartite graph with a left and a right copy of every vertex,
and put an edge $u_Lv_R$ for every arc $u\to v$. It is $k$-regular.
For any subset $X$ of the left side, counting incident edges gives
$k|X|\le k|N(X)|$, so Hall's condition holds. A perfect matching
selects exactly one incoming and one outgoing arc per original
vertex; these arcs are a spanning union of directed cycles.
Delete them and repeat in the $(k-1)$-regular bipartite remainder.
The result is a partition into $k$ spanning cycle factors, proving
at least $k$ arc-disjoint cycles without any connectivity assumption.

This explains one use of exact two-sided regularity. An outregular
but indegree-unbalanced digraph need not have the bipartite matching
condition, so the factorization would not follow from outdegrees
alone. However, a factor can consist of one Hamilton cycle, so this
argument by itself yields only a linear count.

\paragraph{3. A shorter first cycle improves the count to $k+1$.}
Assume $k\ge2$. There is a directed cycle not spanning all $n$
vertices. To prove this, take any cycle. If it is not Hamiltonian,
we are done. Otherwise label its cyclic order $v_0,\ldots,v_{n-1}$.
Some outgoing arc $v_i\to v_j$ differs from its successor arc,
since outdegree is at least two. It is not a loop. This extra arc,
followed by the Hamilton-cycle segment from $v_j$ to $v_i$, forms
a cycle of length at most $n-1$.

Delete this shorter cycle $C$. The remainder is balanced and has
\[
                 kn-|C|>(k-1)n
\]
arcs. Decompose it into cycles by the first argument. Each cycle
has at most $n$ arcs, so the remainder contains at least $k$ cycles.
Together with $C$ there are at least $k+1$ arc-disjoint cycles.
For $k=2$ this is exactly the required three, while $k=1$ follows
from the initial decomposition. For larger $k$ it falls short of
$k(k+1)/2$ and is not competitive with [S3].

\paragraph{4. A full symmetric subclass and an explicit packing.}
In a symmetric $k$-regular digraph, taking every opposite arc pair
as a digon gives $nk/2$ arc-disjoint cycles. Since simplicity implies
$n\ge k+1$, this number is at least $k(k+1)/2$, proving the full
target for this subclass. For a nonsymmetric degree-two example,
the cyclic tournament on five vertices with steps one and two has
the arc-partition into cycles
\[
            (0,1,3,0),\quad(0,2,4,0),\quad(1,2,3,4,1).
\]
These use respectively three, three and four of its ten arcs.

\paragraph{Verification, gap and final status.}
The script [S9] verifies this explicit packing and the symmetric
counts, including absence of duplicated arcs. The first general gap
is forcing sufficiently many short cycles repeatedly; one shorter
cycle yields only one extra member. The matching factorization
provides no such length control. Full target: UNRESOLVED here;
$k=1,2$ and the symmetric subclass are proved independently.
\subsection{Sources}
\begin{itemize}
\item[S1] ULAM.AI and the UnsolvedMath Contributors, supplied problems.json, record 3271 (OPG-49573), OpenGarden collection. Catalogue identity and source transcription; CC BY 4.0.
\url{https://huggingface.co/datasets/ulamai/UnsolvedMath}.
\item[S2] Open Problem Garden, Arc-disjoint directed cycles in regular directed graphs. Original problem and discussion; the mathematical formulation below is expanded from the supplied source record.
\url{http://www.openproblemgarden.org/op/arc_disjoint_directed_cycles_in_regular_directed_graphs}.
\item[S3] Noga Alon, Colin McDiarmid and Michael Molloy, \emph{Edge-Disjoint Cycles in Regular Directed Graphs}, Journal of Graph Theory 22 (1996), 231--237. Author-hosted paper: \url{https://web.math.princeton.edu/~nalon/PDFS/cycles6.pdf}.
\item[S9] GPT-6 Astra Ultra, exact finite checks accompanying this attempt (2 October 2026). Repository script, Python standard library only. \texttt{scripts/check\_attempt\_batch03\_digraphs\_2026\_10\_02.py}.
\end{itemize}
\end{document}
